% Abhakseite „Das kann ich“ – Zone nur im Gesamt (\mitzone)
%% TODO Ich-kann-Sätze: Platzhalter wie in den Aufgabendateien
\begin{abhakseite}
\ifdefined\mitzone
\abhakgruppe{Kennst du schon}
\abhak{1}{Terme mit einem Buchstaben mehr umformen}
\abhak{2}{Lineare Gleichungen nach einem beliebigen Buchstaben auflösen}
\abhak{3}{Quadratische Gleichungen und Diskriminante}
\abhak{4}{Parabelparameter kennen}
\abhak{5}{Ableitungsregeln sicher anwenden, auch mit Vorfaktoren}
\abhak{6}{Extrem- und Wendepunktkriterien an fester Funktion führen}
\abhak{7}{Stammfunktion bilden und bestimmtes Integral auswerten}
\abhak{8}{Terme mit einem Buchstaben mehr umformen – Fehler finden}
\abhak{9}{Terme mit einem Buchstaben mehr umformen – selbst rechnen}
\fi
\abhakgruppe{Einheit 1 · Scharbegriff und Parameterwert}
\abhak{10}{Scharbegriff und Parameterwert}
\abhak{11}{Scharbegriff und Parameterwert – Fehler finden, Begründen, Anwendung}
\abhakgruppe{Einheit 2 · Eigenschaften aller Graphen am Term}
\abhak{12}{Eigenschaften aller Graphen}
\abhak{13}{Eigenschaften aller Graphen – Fehler finden, Begründen}
\abhakgruppe{Einheit 3 · Extrem- und Wendepunkte mit Parameter}
\abhak{14}{Extrem- und Wendepunkte mit Parameter}
\abhak{15}{Extrem- und Wendepunkte mit Parameter – Fehler finden, Begründen, Anwendung}
\abhakgruppe{Einheit 4 · Parameter aus Bedingungen bestimmen}
\abhak{16}{Parameter aus Bedingungen}
\abhak{17}{Parameter aus Bedingungen – Fehler finden, Begründen, Anwendung}
\abhakgruppe{Einheit 5 · Ortskurve und Kurvenvergleich}
\abhak{18}{Ortskurve und Kurvenvergleich}
\abhak{19}{Ortskurve und Kurvenvergleich – Fehler finden, Begründen}
\end{abhakseite}
